\documentclass[tikz, border=10pt]{standalone}
\usepackage{tikz}
\usepackage{amsmath}
\usetikzlibrary{calc, arrows.meta}

\begin{document}
\begin{tikzpicture}[font=\sffamily, scale=2.6]

    % K = 3 case: the simplex Delta^2 embeds in R^3 as the triangle
    % with vertices e_0, e_1, e_2. Drawn here in the 2D centered
    % tangent plane V (a (K-1)-dim subspace of R^K). The radius-r
    % sphere in V is, for K=3, a circle around the reference probability
    % vector p; sample perturbations eps sit on that circle.
    \def\rad{1.2}
    \def\rball{0.34}

    \pgfmathsetmacro{\axc}{\rad*cos(90)}
    \pgfmathsetmacro{\ayc}{\rad*sin(90)}
    \pgfmathsetmacro{\bxc}{\rad*cos(210)}
    \pgfmathsetmacro{\byc}{\rad*sin(210)}
    \pgfmathsetmacro{\cxc}{\rad*cos(330)}
    \pgfmathsetmacro{\cyc}{\rad*sin(330)}

    % Simplex triangle Delta^{K-1}.
    \draw[gray!70, thick]
        (\axc,\ayc) -- (\bxc,\byc) -- (\cxc,\cyc) -- cycle;
    \node[font=\footnotesize, above=1pt] at (\axc,\ayc) {$\mathbf{e}_0$};
    \node[font=\footnotesize, below left=-1pt and 1pt]
        at (\bxc,\byc) {$\mathbf{e}_1$};
    \node[font=\footnotesize, below right=-1pt and 1pt]
        at (\cxc,\cyc) {$\mathbf{e}_2$};

    % Reference probability vector p in the interior, slightly off-center
    % (matching the typical Dirichlet(alpha=2) sample).
    \def\pX{0.18}
    \def\pY{0.05}
    \fill[black] (\pX,\pY) circle (0.018);
    \node[font=\footnotesize, anchor=north west]
        at (\pX+0.03,\pY+0.02) {$\mathbf{p}$};

    % Radius-r sphere in V: at K=3, V is 2-dim and the sphere is a
    % circle around p. Drawn dotted to suggest "set of allowed
    % perturbations".
    \draw[gray!55, dotted, thick] (\pX,\pY) circle (\rball);

    % A handful of sample eps perturbations on the sphere (arrows from p
    % to the circle). Angles chosen to avoid overlap with the labels.
    \foreach \ang in {30, 110, 200, 285} {
        \pgfmathsetmacro{\ux}{\pX + \rball*cos(\ang)}
        \pgfmathsetmacro{\uy}{\pY + \rball*sin(\ang)}
        \draw[->, gray!60!black, thick, >=Stealth]
            (\pX,\pY) -- (\ux,\uy);
        \fill[gray!60!black] (\ux,\uy) circle (0.012);
    }

    % Annotation for the radius.
    \pgfmathsetmacro{\rax}{\pX + \rball*cos(285)*0.55}
    \pgfmathsetmacro{\ray}{\pY + \rball*sin(285)*0.55}
    \node[font=\scriptsize, color=gray!55!black, anchor=west]
        at ($(\pX,\pY)+(0.18,-0.16)$) {$r$};

    % Legend: a table centred under the simplex, so the drawing stays
    % centred and the entries cannot overlap. The inline symbols use
    % scale=1 so they do not inherit the picture's scale.
    \node[anchor=north, font=\footnotesize, inner sep=0pt]
        at (0, \byc-0.22) {%
        \begin{tabular}{@{}c@{\hspace{0.5em}}l@{\hspace{1.6em}}c@{\hspace{0.5em}}l@{}}
            \tikz[baseline=-0.6ex, scale=1]
                \draw[gray!70, thick] (0, 0) -- (0.45, 0);
            & simplex $\Delta^{K-1}$
            & \tikz[baseline=-0.6ex, scale=1]
                \fill[black] (0, 0) circle (1.2pt);
            & reference probability vector $\mathbf{p}$ \\[2pt]
            \tikz[baseline=-0.6ex, scale=1]
                \draw[gray!55, dotted, thick] (0, 0) circle (4pt);
            & radius-$r$ sphere in $V$
            & \tikz[baseline=-0.6ex, scale=1]
                \draw[->, gray!60!black, thick, >=Stealth] (0, 0) -- (0.45, 0);
            & sample perturbation $\boldsymbol{\epsilon}\sim\mathrm{Unif}(\partial B_r)$
        \end{tabular}};

\end{tikzpicture}
\end{document}
